\scanpage{30}
\begin{center}6 - 41 - \hfill n° 392

\S 5.\end{center}
1) Soient $A$ un anneau noethérien intégralement clos, $K$ son corps des fractions, $L$ une extension de degré $n$ de $K$, $B$ la fermeture intégrale de $A$ dans $L$ ; on suppose que $B$ est un $A$-module de type fini. Soient $E$, $F$ deux anneaux tels que $A\subset E\subset F\subset B$.

a) Montrer que pour tout idéal premier $\uline{p}$ de hauteur 1 dans $E$, il n'y a qu'un nombre fini d'idéaux premiers $\uline{P}$ de $F$ au-dessus de $\uline{p}$ et qu'ils sont de hauteur 1 ; réciproquement, pour tout idéal premier $\uline{P}$ de hauteur 1 dans $F$, $\uline{P}\cap E$ est de hauteur 1 dans $E$.

b) Pour tout idéal premier $\uline{P}$ de hauteur 1 dans $F$, on appelle \uline{degré} de $\uline{P}$ par rapport à $E$, et on note $\deg_E(\uline{P})$ le degré du corps des fractions de $F/\uline{P}$ sur le corps des fractions de $E/(\uline{P}\cap E)$ ; le degré par rapport à $E$ d'un idéal premier de $F$ qui n'est pas de hauteur 1 est pris égal à 0 par définition. Si $\uline{Q}$ est un idéal primaire de $F$, $\uline{P}$ l'idéal premier correspondant, et si $m$ est le plus petit entier tel que la puissance symbolique $\uline{P}^{(m)}$ de $\uline{P}$ soit contenue dans $\uline{Q}$, on pose $\deg_E(\uline{Q})=m.\deg_E(\uline{P})$, (de sorte que $\deg_F(\uline{Q})=m$ si $\uline{P}$ est de hauteur 1, $\deg_F(\uline{Q})=0$ dans le cas contraire). Finalement, si $\uline{a}$ est un idéal quelconque $\ne0$ de $F$, $\uline{Q}_i$ les idéaux primaires de la décomposition primaire réduite de $\uline{a}$, correspondant aux idéaux premiers non immergés, on pose $\deg_E(\uline{a})=\sum_i\deg_E(\uline{Q}_i)$.

Avec ces notations, soit $C$ un anneau tel que $A\subset C\subset B$ ; on dit qu'un idéal $\uline{a}$ de $C$ est \uline{pseudo-inversible} si, pour tout idéal premier $\uline{p}$ de hauteur 1 dans $A$, $\uline{a}_{\uline{p}}$ est inversible dans l'anneau semi-local $C_{\uline{p}}$. Montrer que, pour que $\uline{a}\ne0$ soit pseudo-inversible, il faut et il suffit que l'on ait $\deg_C(\uline{a})=\sum_{\uline{P}}n(\uline{P})\deg_C(\uline{P})$, où $\uline{P}$ parcourt l'ensemble des idéaux premiers de hauteur 1 de $B$ et les $n(\uline{P})$ sont définis par $\operatorname{div}(\uline{a}B)=\sum_{\uline{P}}n(\uline{P})\uline{P}$. (Se ramener au cas où $A$ est un anneau de valuation discrète, donc $C$ semi-local. Observer d'une part que dans $C$, il y a alors identité entre les idéaux inversibles et les idéaux principaux, et que si $\uline{a}=xC$, $\deg_C(\uline{a})$ est égal à la longueur du $C$-module $C/xC$, et aussi à celle du $C$-module $B/xB$. D'autre part, si $\uline{m}_i$, $(1\le i\le r)$, sont les idéaux maximaux distincts de $C$, $\uline{a}\ne0$ un idéal de $C$, il existe $y\in\uline{a}$ tel que $y\notin C\cap(\uline{a}\uline{m}_iB)$ pour tout $i$ ; en déduire que $yB=\uline{a}B$, en utilisant le fait que $B$ est alors un anneau principal).
Si $C=B$, tout idéal $\ne0$ de $B$ est pseudo-inversible.
\mspage{29}
